%! Two Categorical Variables — Unit 2, Lesson 2.1 — AP Practice
% GRADUAL RELEASE: AP-format practice, assigned but NOT scored — the cover's score
% column reads NA for this component. The graded practice is the `homework` component that
% follows it at the back of the packet.
% ONE PAGE (compact budget): two MC items + one FRQ of three parts, on one grocery survey.
% MC 1 carries the denominator trap (distractor B is the reverse conditional, 60/250).
% MC 2 carries the crux (distractor A compares raw counts from groups of 250 and 150).
% FRQ (c) is the spiral to Lessons 1.2–1.3 (a bar graph is for categories, a histogram is not).
% Page lockstep with ap_practice_key rests on \writelines{n} in the blank matching exactly n
% terse \ansline{} calls in the key.
% The next-lesson preview is NOT here: it closes the packet, in ../homework/.
%
% DATA (400 households, location × main way of getting groceries):
%            In store  Delivery  Pickup  Total
%   Urban       150       60       40     250   -> 60% / 24% / 16%
%   Rural       105       15       30     150   -> 70% / 10% / 20%
%   Total       255       75       70     400
\documentclass[12pt]{article}
\usepackage{apstats-article}
\usepackage{apstats-boxes}

\begin{document}

\pageheader{Unit 2, Lesson 2.1}{AP Practice}
% NAMESTRIP: no \namedateperiod here — the name row lives on the cover only.

\vspace{0.04in}

% Authored BYTE-IDENTICALLY in ap_practice/ and ap_practice_key/ — this box is what keeps the
% blank and the key the same number of pages.
\boxguard[8]
\begin{remindbox}
\small
This page is \textbf{not required and not scored} --- your graded homework is the last section
of this packet. Work these AP-format reps and bring what stalls you to study hall.
\end{remindbox}

\vspace{0.05in}

\boxguard[20]
\begin{scenariobox}[Getting the Groceries]{navy}
\small
\begin{minipage}[c]{0.30\linewidth}\raggedright
A grocery chain surveyed 400 households: where they live, and how they \emph{usually} get
groceries.
\end{minipage}\hfill
\begin{minipage}[c]{0.68\linewidth}
\centering
\renewcommand{\arraystretch}{1.1}
\begin{tabular}{@{} l | c c c | c @{}}
 & \textbf{In store} & \textbf{Delivery} & \textbf{Curbside pickup} & \textbf{Total} \\ \hline
\textbf{Urban} & 150 & 60 & 40 & 250 \\
\textbf{Rural} & 105 & 15 & 30 & 150 \\ \hline
\textbf{Total} & 255 & 75 & 70 & 400 \\
\end{tabular}
\end{minipage}
\end{scenariobox}

\vspace{0.05in}

\boxguard[12]
\begin{headlinebox}{goldbg}
\small
\textbf{Section I --- Multiple Choice.} Circle the single best answer.
\end{headlinebox}

\vspace{0.04in}

\begin{enumerate}[leftmargin=1.9em, itemsep=0.08in]

  \item Of the households that use delivery, what proportion are urban?
        \par\vspace{3pt}
        \hspace*{1.2em}\textbf{(A)} 0.150 \hfill \textbf{(B)} 0.240 \hfill \textbf{(C)} 0.625
        \hfill \textbf{(D)} 0.800 \hfill \textbf{(E)} 0.188 \hspace*{1.2em}

  \item Which statement is best supported by the survey?
        \begin{enumerate}[label=\textbf{(\Alph*)}, leftmargin=2.2em, itemsep=0pt]
          \item Urban households shop in store more than rural ones, since 150 exceeds 105.
          \item Location and grocery method are associated: 24\% of urban households use
                delivery, vs.\ 10\% of rural households.
          \item There is no association: in both locations, most households shop in store.
          \item Living in a city causes a household to order grocery delivery.
          \item The variables cannot be associated, because both are categorical.
        \end{enumerate}

\end{enumerate}

\vspace{0.06in}

\boxguard[12]
\begin{headlinebox}{goldbg}
\small
\textbf{Section II --- Free Response.} Show your work, and answer \emph{in context}.
\end{headlinebox}

\vspace{0.04in}

\begin{enumerate}[leftmargin=1.9em, itemsep=0.08in, resume]

  \item \textbf{(a)} Find the conditional distribution of grocery method for \emph{rural}
        households, as percents.\\[4pt]
        \writelines{1}

        \vspace{4pt}
        \textbf{(b)} Is there an association between location and grocery method? Justify.\\[4pt]
        \writelines{2}

        \vspace{4pt}
        \textbf{(c)} Why is a segmented bar graph, not a histogram, the right display here?\\[4pt]
        \writelines{2}

\end{enumerate}

\end{document}
